\section*{Annexes}
\addcontentsline{toc}{section}{Annexes}
\renewcommand{\thesubsection}{\Alph{subsection}}
\setcounter{subsection}{0}
\renewcommand{\theequation}{\Alph{subsection}.\arabic{equation}}

\subsection{Moments gaussiens du modèle}
\label{ann:moments}
\setcounter{equation}{0}

Notons $(z_1, z_2) = (x, y)$, de paramètres $(a_1, \sigma_1) = (a_x, \sigma_x)$ et $(a_2, \sigma_2) = (a_y, \sigma_y)$,
avec $\rho_{12} = \rho$ et $\rho_{11} = \rho_{22} = 1$. Avec $\tau = T - t$ et $H_a = H_a(t,T)$, les moments
conditionnels à $\F_t$ utilisés aux sections~\ref{sec:modele} à~\ref{sec:resultats} sont, pour $i, j \in \{1, 2\}$ :
\begin{align}
\Cov\big(z_i(T), z_j(T)\big) &= \rho_{ij}\sigma_i\sigma_j\,\frac{1-e^{-(a_i+a_j)\tau}}{a_i+a_j}, \\
\Cov\Big(z_i(T), \int_t^T z_j\Big) &= \frac{\rho_{ij}\sigma_i\sigma_j}{a_j}\Big[\frac{1-e^{-a_i\tau}}{a_i} - \frac{1-e^{-(a_i+a_j)\tau}}{a_i+a_j}\Big], \\
\Cov\Big(\int_t^T z_i, \int_t^T z_j\Big) &= \frac{\rho_{ij}\sigma_i\sigma_j}{a_ia_j}\Big[\tau - H_{a_i} - H_{a_j} + \frac{1-e^{-(a_i+a_j)\tau}}{a_i+a_j}\Big].
\end{align}
Pour $i = j$, la troisième expression redonne $V_{a,\sigma}$, et, pour $i \neq j$, deux fois cette covariance est le
terme croisé $V_{xy}$. Les écarts-types et la corrélation de $(x(T_0), y(T_0))$ qui entrent dans le prix
exact~\eqref{eq:prix_exact} sont $s_x^2 = \Var x(T_0)$, $s_y^2 = \Var y(T_0)$ et
$r_{xy} = \Cov(x(T_0), y(T_0))/(s_xs_y)$, calculés en $t = 0$.

\paragraph{Moyennes sous la mesure forward.} D'après~\eqref{eq:zc_hw}, le numéraire $P(u,T_0)$ a pour
volatilité $-\sigma_x H_{a_x}(u,T_0)$, portée par le seul brownien $W_x$. Le changement de numéraire donne
\begin{equation}
\dd W_x^{T_0} = \dd W_x + \sigma_x H_{a_x}(u,T_0)\,\dd u,\qquad
\dd W_y^{T_0} = \dd W_y + \rho\,\sigma_x H_{a_x}(u,T_0)\,\dd u.
\end{equation}
En intégrant la dynamique~\eqref{eq:dynamique} de $0$ à $T_0$, il vient
\begin{equation}
\begin{split}
\mu_x &= -\int_0^{T_0}\sigma_x^2 H_{a_x}(u,T_0)\,e^{-a_x(T_0-u)}\,\dd u = -\Cov\Big(x(T_0),\int_0^{T_0}x\Big),\\
\mu_y &= -\int_0^{T_0}\rho\,\sigma_x\sigma_y H_{a_x}(u,T_0)\,e^{-a_y(T_0-u)}\,\dd u = -\Cov\Big(y(T_0),\int_0^{T_0}x\Big),
\end{split}
\end{equation}
dont les expressions fermées sont celles de la section~\ref{sec:evaluation}.

\subsection{Prix des options}
\label{ann:prix}

\begin{table}[H]
\centering\small
\caption{Calls et puts sur l'OAT 4,10~\% 2046, modèle à deux facteurs, prix exacts pour 100 de nominal (strike
clean en \% du forward clean)}
\label{tab:prix_complets}
\begin{tabular}{@{}lrrrrr@{}}
\toprule
Strike & 1 an & 2 ans & 5 ans & 7 ans & 10 ans \\
\midrule
\multicolumn{6}{@{}l}{\emph{Calls}} \\
90~\%  & 9,434 & 9,837 & 9,961 & 9,514 & 8,442 \\
95~\%  & 6,099 & 6,951 & 7,660 & 7,421 & 6,539 \\
100~\% & 3,578 & 4,677 & 5,763 & 5,679 & 4,960 \\
105~\% & 1,899 & 3,000 & 4,248 & 4,271 & 3,689 \\
110~\% & 0,914 & 1,840 & 3,073 & 3,161 & 2,695 \\
\midrule
\multicolumn{6}{@{}l}{\emph{Puts}} \\
90~\%  & 0,641 & 1,398 & 2,475 & 2,577 & 2,206 \\
95~\%  & 1,703 & 2,732 & 3,917 & 3,952 & 3,421 \\
100~\% & 3,578 & 4,677 & 5,763 & 5,679 & 4,960 \\
105~\% & 6,296 & 7,220 & 7,990 & 7,740 & 6,807 \\
110~\% & 9,707 & 10,279 & 10,558 & 10,099 & 8,931 \\
\midrule
\multicolumn{6}{@{}l}{\emph{Forward et espérance du modèle (prix dirty)}} \\
$F$ (repo) & 91,573 & 90,525 & 87,908 & 86,809 & 86,445 \\
$M$ (non recentré) & 91,267 & 90,172 & 87,824 & 87,516 & 88,742 \\
\bottomrule
\end{tabular}
\end{table}

\subsection{Ajustement de la courbe OAT}
\label{ann:oat}

\begin{table}[H]
\centering\small
\caption{OAT utilisées pour la courbe Nelson-Siegel-Svensson : prix clean au 8 septembre 2026 et résidus en
rendement}
\label{tab:residus_oat}
\begin{tabular}{@{}llccr@{}}
\toprule
ISIN & Maturité & Coupon (\%) & Prix clean & Résidu (pb) \\
\midrule
FR0129570570 & 30/09/2026 & 0,00 & 99,868 & $-0{,}4$ \\
FR0129704112 & 09/12/2026 & 0,00 & 99,372 & 1,4 \\
FR0129704153 & 10/03/2027 & 0,00 & 98,649 & $-1{,}5$ \\
FR0129570646 & 16/06/2027 & 0,00 & 97,849 & 0,7 \\
FR0129704195 & 08/09/2027 & 0,00 & 97,121 & $-0{,}8$ \\
FR001400XLW2 & 24/09/2028 & 2,40 & 98,506 & 0,7 \\
FR0014016G71 & 24/09/2029 & 2,40 & 97,474 & 0,5 \\
FR0013516549 & 25/11/2030 & 0,00 & 86,758 & 1,3 \\
FR0014018OI0 & 25/02/2032 & 3,25 & 98,160 & $-1{,}7$ \\
FR001400BKZ3 & 25/11/2032 & 2,00 & 90,604 & 0,4 \\
FR001400L834 & 25/11/2033 & 3,50 & 97,730 & $-1{,}2$ \\
FR001400QMF9 & 25/11/2034 & 3,00 & 93,115 & $-0{,}6$ \\
FR0014012II5 & 25/11/2035 & 3,50 & 95,405 & 0,1 \\
FR0014018YR0 & 25/11/2036 & 3,70 & 95,759 & 0,0 \\
FR001400WYO4 & 25/05/2042 & 3,60 & 88,794 & 2,8 \\
\textbf{FR0014015MU5} & \textbf{25/05/2046} & \textbf{4,10} & \textbf{91,690} & \textbf{2,0} \\
FR0013404969 & 25/05/2050 & 1,50 & 53,283 & 8,4 \\
FR0014016CV2 & 25/05/2057 & 4,40 & 90,960 & $-14{,}5$ \\
FR0014001NN8 & 25/05/2072 & 0,50 & 22,699 & 3,2 \\
\midrule
\multicolumn{5}{@{}l}{\footnotesize Erreur quadratique moyenne : 4,1~pb. En gras, le sous-jacent des options.} \\
\bottomrule
\end{tabular}
\end{table}

\subsection{Environnement de calcul}
\label{ann:calcul}

Les calculs sont regroupés dans un notebook Jupyter unique, \nolinkurl{main/oat_option_pricer.ipynb}, écrit en
Python~3.12 ; l'environnement est figé par le gestionnaire \texttt{uv} (fichiers \texttt{pyproject.toml} et
\texttt{uv.lock}). Les bibliothèques utilisées sont NumPy et SciPy (optimisation, quadrature de Gauss-Hermite,
splines), pandas, matplotlib, statsmodels (filtre de Kalman par la classe \texttt{MLEModel}, régression à
changement de régime \texttt{MarkovRegression}, tests ARCH-LM et de Ljung-Box), scikit-learn (analyse en
composantes principales) et joblib (parallélisation du test sur données simulées). L'exécution complète du
notebook prend environ une minute et quinze secondes. Les données de marché sont lues dans le fichier Bloomberg
\nolinkurl{data/market data bloom.xlsx} et la courbe repo fictive dans \nolinkurl{data/repo_fictif.csv} ; une courbe de
marché s'y substituerait au même format. Les figures de ce rapport sont produites par le script
\nolinkurl{memoire/figures/generer_figures.py}, qui exécute le notebook sans modifier ses calculs.
